\section{Exact obstructions for two specified trace representations}
\label{sec:per-trace-obstructions}

Retain the parameter projection, source fibre, and multiplication operator
of the preceding section. Write $\mathcal A_N$ for its full split source
algebra and $u_Z$ for the projected multiplier, where
$N=3n^2$, $M=n^3+6n^2$. In its original point-idempotent basis the
diagonal entries of multiplication by $u_Z$ are zero except at the
uniquely flagged permutation witnesses. Their entries are
$\prod_i Z_{i,\sigma(i)}$. Thus the already constructed polynomial is
\[
 P_n(Z)=\operatorname{Tr}_{\mathcal A_N}(m_{u_Z})
       =\sum_{\sigma\in S_n}\prod_{i=1}^n Z_{i,\sigma(i)}.
\]
The following calculations concern this same trace. No change of source
coordinates or removal of its zero eigenvalues is needed. The original
rational-coordinate fibre retains its characteristic-zero base here.
The integral split idempotent lattice constructed above carries the
same diagonal multiplication matrix with entries in $\mathbb Z[Z]$;
its trace polynomial therefore has integer coefficients. It is that
integral matrix and polynomial, and hence their coefficient matrices,
which can be base changed to any field, including characteristic two.
This does not reduce the source coordinates containing $1/2$ modulo two.

\subsection{Coefficient contraction and its exact rank}

Fix a field $k$, $n\ge1$, and $0\le k_0\le n$. To avoid using the
field symbol as an integer index, denote the cut by $k_0$ throughout.
Let $H_{k_0}$ be the vector space with basis
\[
 h_{\boldsymbol j}=\prod_{i=1}^{k_0}Z_{i,j_i},
 \qquad \boldsymbol j\in\{1,\ldots,n\}^{k_0},
\]
and let $T_{k_0}$ have basis
$t_{\boldsymbol l}=\prod_{i=k_0+1}^{n}Z_{i,l_i}$ indexed by the
remaining $n-k_0$ choices. Empty products give the one-dimensional
spaces at the two endpoints. Since the variable sets are disjoint,
multiplication identifies $H_{k_0}\otimes_k T_{k_0}$ with the space
of polynomials homogeneous of degree one in each row of $Z$.
Indeed, multiplication maps the displayed tensor basis bijectively
onto its distinct monomial basis. Define the coefficient contraction
\[
 \mathcal F_{k_0}:H_{k_0}^*\longrightarrow T_{k_0},\qquad
 \mathcal F_{k_0}(h_{\boldsymbol j}^*)
  =\sum_{\boldsymbol l}
   [h_{\boldsymbol j}t_{\boldsymbol l}]P_n\,t_{\boldsymbol l}.
\]
Here $[m]$ denotes the coefficient of the exact monomial $m$.

\begin{theorem}\label{thm:per-exact-row-rank}
For every field and every displayed cut,
$\operatorname{rank}\mathcal F_{k_0}=\binom n{k_0}$.
There is an integral factorization through the free module on the
$k_0$-element subsets of $\{1,\ldots,n\}$, and an integral identity
minor of that size.
\end{theorem}
\begin{proof}
Let $U_{k_0}$ have basis $e_S$ for those subsets. Define
$a:H_{k_0}^*\to U_{k_0}$ by
$a(h_{\boldsymbol j}^*)=e_{\{j_1,\ldots,j_{k_0}\}}$ if the
indices $j_i$ are distinct, and by zero otherwise. Define
\[
 b(e_S)=\sum_{\{l_{k_0+1},\ldots,l_n\}=S^c
                   \atop\text{all indices distinct}}t_{\boldsymbol l}.
\]
A monomial in $P_n$ uses each column exactly once. Conversely every
choice using each row and column exactly once is the monomial of a
unique permutation and has coefficient one. These two facts give
$\mathcal F_{k_0}=ba$ on every displayed basis vector, including
those with repeated columns. The factorization proves the upper bound.

For each $S$, take $\boldsymbol j(S)$ to be the elements of $S$ in
increasing order assigned to the first $k_0$ rows, and take
$\boldsymbol l(S)$ to be the elements of $S^c$ in increasing order
assigned to the remaining rows. The coefficient of
$h_{\boldsymbol j(S)}t_{\boldsymbol l(S')}$ is one exactly when
$S=S'$: disjointness and cardinalities force $S=(S'^c)^c$.
It is zero otherwise. Thus these rows and columns give the identity
matrix over $\mathbb Z$, proving the lower bound over every field.
No division by a factorial or by a characteristic-dependent integer
has occurred.
\end{proof}

This is also an exact map from the original trace construction.
For a matrix $B(Z)$ over a polynomial ring, finite coefficient extraction
commutes with its trace:
$[m]\operatorname{Tr}B=\sum_v[m]B_{vv}
=\operatorname{Tr}([m]B)$. This follows directly from the definition
of the diagonal sum, for every coefficient and over every base ring.
Apply it to the full $3^N\times3^N$ matrix $m_{u_Z}$.
The projection to its uniquely flagged permutation diagonal entries
therefore gives precisely the maps $a,b$ above, while every retained
zero diagonal entry contributes zero. This proves the coefficient
contraction correspondence; the matrix size $3^N$ is unchanged.

For every integer $d\ge1$ the same argument applied to the previously
proved power trace $\operatorname{Tr}(m_{u_Z}^d)
=\operatorname{per}_n((Z_{ij}^d))$ gives the same rank.
Use the head and tail bases with each exponent $1$ replaced by $d$;
these are still distinct monomials in every characteristic. Their
coefficient matrix is the same integral zero-one matrix.

\subsection{Optimal widths of row-ordered branching programs}

A row-ordered algebraic branching program here means a finite directed
layered graph with layers $0,1,\ldots,n$, one source at layer zero,
one sink at layer $n$, and edges only from layer $i-1$ to layer $i$.
Every such edge is labelled by a homogeneous linear form in
$Z_{i1},\ldots,Z_{in}$ over $k$. The output is the sum over all
source-to-sink paths of the product of their labels. The width $w_i$
is the number of vertices at layer $i$; distinct parallel edges may
equivalently be combined by adding their labels.

\begin{theorem}\label{thm:per-row-abp}
Every such program for $P_n$ has $w_i\ge\binom ni$ at every layer.
These bounds hold simultaneously with equality. The minimum total
number of vertices, including both endpoints, is $2^n$.
An attaining program has exactly $n2^{n-1}$ edges.
\end{theorem}
\begin{proof}
For a vertex $v$ in layer $i$, let $p_v$ be the sum of path products
from the source to $v$, and $q_v$ the corresponding sum from $v$
to the sink. Each full path has exactly one such $v$ and decomposes
uniquely into its prefix and suffix. Distributivity therefore gives
$P_n=\sum_{v\text{ in layer }i}p_vq_v$, with $p_v\in H_i$ and
$q_v\in T_i$. Its coefficient contraction factors as
\[
 H_i^*\longrightarrow k^{w_i}\longrightarrow T_i,
 \qquad\lambda\longmapsto(\lambda(p_v))_v,
 \qquad(c_v)_v\longmapsto\sum_v c_vq_v.
\]
Its rank is at most $w_i$. The preceding theorem proves the lower bound.

For attainment, put at layer $i$ all subsets $S\subseteq\{1,\ldots,n\}$
of size $i$. For every $j\notin S$ put an edge from $S$ to
$S\cup\{j\}$ labelled $Z_{i+1,j}$. A full path chooses the columns
$j_1,\ldots,j_n$ without repetition, so determines one permutation.
Every permutation determines exactly this path. Multiplying its
labels and then summing gives $P_n$. The number of vertices is
$\sum_i\binom ni=2^n$, because each subset has one cardinality and
there are two membership choices for each of the $n$ elements.
The number of edges is
$\sum_{S\subseteq[n]}(n-|S|)=n2^{n-1}$: count pairs $(S,j)$ with
$j\notin S$, choosing $j$ in $n$ ways and an arbitrary subset of
its $n-1$ other elements. Summing the width lower bounds proves
minimality of the stated total vertex number.
\end{proof}

The row condition is part of this precise representation class.
The factorization just proved does not apply to a branching program
that may reuse a row at several layers or intermix row variables
inside a layer. Thus the result gives no lower bound for that larger
class from an unproved identification of its states with $H_i$.
The proved exact map is the prefix--suffix contraction above.

\subsection{An exponential obstruction for nonnegative arithmetic}

Now work over $\mathbb R$. A monotone arithmetic circuit means a
finite acyclic circuit with leaves that are variables $Z_{ij}$ or
nonnegative real constants, and binary addition and multiplication
gates. Its size $s$ is the number of arithmetic-operation gates.
Fan-out is allowed. No subtraction or division is allowed. A parse
tree at a gate is defined recursively: at an addition choose one
of its two children; at a multiplication take both; copy occurrences
when a gate is used more than once. A nonzero parse tree is one
whose constant-leaf product is positive. It contributes that positive
coefficient times the product of its variable leaves.
Induction in a topological ordering proves that a gate polynomial
is the sum of its parse-tree contributions. There are only finitely
many parse trees, since the circuit is finite and acyclic.

\begin{lemma}\label{lem:per-live-gate-profile}
Suppose such a circuit computes $P_n$. A gate occurrence is called
live if it occurs in a nonzero output parse tree, and a gate is live
if it has such an occurrence. For every live gate $g$ there are
uniquely determined subsets $I_g,J_g\subseteq[n]$ of the same size
such that every monomial at $g$ with positive coefficient is a
matching from $I_g$ onto $J_g$. In particular its degree is the
fixed integer $d_g=|I_g|=|J_g|$.
\end{lemma}
\begin{proof}
Fix one nonzero output parse tree and one occurrence of $g$ in it.
Remove the subtree rooted at that occurrence and retain the rest,
including every other occurrence of $g$ if there is one. Let $c m$
be the product of the retained constant and variable leaves, so
$c>0$. Every monomial $t$ with positive coefficient at $g$ has a
nonzero parse tree there, by the parse expansion. Substituting that
tree in the retained occurrence produces a nonzero output parse
tree contributing a positive multiple of $mt$. Nonnegative
coefficients cannot cancel, so $mt$ occurs in $P_n$.

Every monomial of $P_n$ has total exponent one in each row and
column, and exponent at most one on each variable. Consequently
$m$ is a partial matching; all such $t$ use exactly the rows and
columns missing from $m$, once each. These two complementary
sets are $I_g$ and $J_g$. They have the same size because any
monomial $t$ uses one variable per used row and column. At least
one $t$ exists, since the fixed occurrence was nonzero. That
monomial already determines its sets of rows and columns,
showing uniqueness and independence of the chosen live occurrence.
\end{proof}

At a live addition, every child chosen in a nonzero parse tree has
the same profile and degree as the parent: each of its monomials
also occurs at the parent. At a live multiplication with nonzero
children, both children are live by extending the fixed context,
their row sets and their column sets are disjoint, and their unions
are those of the parent. This follows by multiplying any positive
monomial of the first child by one of the second and using the
lemma. In particular their degrees add. These conclusions establish
the needed homogeneity and disjointness from positivity; they have
not been added as circuit hypotheses.

\begin{theorem}\label{thm:per-monotone-lower-bound}
For every $n\ge4$, every monotone arithmetic circuit computing
the original projected trace $P_n$ satisfies
\[
 s\ \ge\ \binom n{\lceil n/3\rceil}
       \ \ge\ 2^{\lceil n/3\rceil}.
\]
\end{theorem}
\begin{proof}
Each of the $n!$ permutation monomials has positive coefficient,
so choose one of its nonzero output parse trees, using any fixed
ordering of the finite tree set to make the choice definite.
Start at its output occurrence, of degree $n$. At an addition
follow its chosen child. At a multiplication follow a child of
larger degree, with a fixed rule for ties. Continue until reaching
the first occurrence of degree at most $2n/3$.

This procedure terminates: it follows a descending path in a finite
tree, and leaf degrees are zero or one. An addition cannot be the
first degree decrease across the threshold, because its chosen
child has the same degree. Hence that first crossing follows a
multiplication of degree greater than $2n/3$. Its selected child
has at least half the parent degree. The selected gate $g$ thus has
\[
 n/3<d_g\le2n/3,
 \quad\text{and in particular}\quad
 \lceil n/3\rceil\le d_g\le\lfloor2n/3\rfloor.
\]
For $n\ge4$ this degree is at least two, so $g$ is an
arithmetic-operation gate, not an input leaf.

Assign the original permutation to this gate. Its chosen parse tree
contains a monomial from $g$, which matches precisely $I_g$ to $J_g$.
Therefore every permutation assigned to $g$ sends $I_g$ onto $J_g$.
For $d_g=d$ the number of all permutations with this property is
exactly $d!(n-d)!$: choose independently a bijection $I_g\to J_g$
and one $I_g^c\to J_g^c$. Each has $d!$ and $(n-d)!$ choices,
respectively, by successive choices of images. The assignments
partition the $n!$ permutations among at most $s$ gates. Thus
\[
 n!\le s\max_{\lceil n/3\rceil\le d\le\lfloor2n/3\rfloor}
                         d!(n-d)!.
\]
The ratio $\binom n{d+1}/\binom nd=(n-d)/(d+1)$ proves that
the binomial coefficients increase up to their middle and then
decrease symmetrically, since $\binom nd=\binom n{n-d}$.
The two endpoints of the displayed range are complementary
integers, so its minimum binomial coefficient is
$\binom n{\lceil n/3\rceil}$. This proves the first bound.
Finally set $r_0=\lceil n/3\rceil$. For $n\ge4$ we have
$r_0\le n/2$: write $n=3a+b$ with $b=0,1,2$ and check the
three inequalities directly, with $a\ge1$ and the case $b=0$
having $a\ge2$ if necessary. For $0\le i<r_0$,
$(n-i)/(r_0-i)\ge n/r_0\ge2$, since
$r_0(n-i)-n(r_0-i)=i(n-r_0)\ge0$. Multiplication gives
\[
 \binom n{r_0}
 =\prod_{i=0}^{r_0-1}\frac{n-i}{r_0-i}
 \ge2^{r_0},
\]
as claimed.
\end{proof}

This theorem constrains positive arithmetic representations of the
specific original-fibre trace after the proved permanent projection.
The universal verifier polynomial before that projection contains
signed coefficients, and the support lemma was proved only for the
nonnegative circuit just defined. The projection is nevertheless the
exact gate-preserving input-substitution map of the preceding section.
If a nonnegative circuit description survives that map as a circuit
of this class, the displayed lower bound applies to its arithmetic
gate count. For signed circuits the cancellation step in the proof
is unavailable; the present theorem makes no assertion that it holds.
The row-ordered and nonnegative bounds provide two explicit
asymptotic obstructions with their maps and domains. The broader
unrestricted obstruction in the research programme remains to be
constructed.
